%2multibyte Version: 5.50.0.2953 CodePage: 936

\documentclass{article}
%%%%%%%%%%%%%%%%%%%%%%%%%%%%%%%%%%%%%%%%%%%%%%%%%%%%%%%%%%%%%%%%%%%%%%%%%%%%%%%%%%%%%%%%%%%%%%%%%%%%%%%%%%%%%%%%%%%%%%%%%%%%%%%%%%%%%%%%%%%%%%%%%%%%%%%%%%%%%%%%%%%%%%%%%%%%%%%%%%%%%%%%%%%%%%%%%%%%%%%%%%%%%%%%%%%%%%%%%%%%%%%%%%%%%%%%%%%%%%%%%%%%%%%%%%%%
\usepackage{amsmath}
\usepackage{amsfonts}

\setcounter{MaxMatrixCols}{10}
%TCIDATA{OutputFilter=LATEX.DLL}
%TCIDATA{Version=5.50.0.2953}
%TCIDATA{Codepage=936}
%TCIDATA{<META NAME="SaveForMode" CONTENT="1">}
%TCIDATA{BibliographyScheme=Manual}
%TCIDATA{Created=Monday, November 08, 2021 17:14:57}
%TCIDATA{LastRevised=Wednesday, November 10, 2021 12:36:26}
%TCIDATA{<META NAME="GraphicsSave" CONTENT="32">}
%TCIDATA{<META NAME="DocumentShell" CONTENT="Standard LaTeX\Blank - Standard LaTeX Article">}
%TCIDATA{CSTFile=40 LaTeX article.cst}

\newtheorem{theorem}{Theorem}
\newtheorem{acknowledgement}[theorem]{Acknowledgement}
\newtheorem{algorithm}[theorem]{Algorithm}
\newtheorem{axiom}[theorem]{Axiom}
\newtheorem{case}[theorem]{Case}
\newtheorem{claim}[theorem]{Claim}
\newtheorem{conclusion}[theorem]{Conclusion}
\newtheorem{condition}[theorem]{Condition}
\newtheorem{conjecture}[theorem]{Conjecture}
\newtheorem{corollary}[theorem]{Corollary}
\newtheorem{criterion}[theorem]{Criterion}
\newtheorem{definition}[theorem]{Definition}
\newtheorem{example}[theorem]{Example}
\newtheorem{exercise}[theorem]{Exercise}
\newtheorem{lemma}[theorem]{Lemma}
\newtheorem{notation}[theorem]{Notation}
\newtheorem{problem}[theorem]{Problem}
\newtheorem{proposition}[theorem]{Proposition}
\newtheorem{remark}[theorem]{Remark}
\newtheorem{solution}[theorem]{Solution}
\newtheorem{summary}[theorem]{Summary}
\newenvironment{proof}[1][Proof]{\noindent\textbf{#1.} }{\ \rule{0.5em}{0.5em}}
\input{tcilatex}
\begin{document}


\begin{center}
{\Large Class Test}
\end{center}

EACH QUESTION\ EQUALLY\ MARKED\ AT\ 12.5 MARKS\bigskip

\textbf{Question 1}

(a) Define an efficient estimator and explain why this is an important
property and how it relates to the mean squared error of an estimator.

Solution:

An estimator $\hat{\theta}_{n}$ is called efficient if it is unbiased and
has the smallest variance, that is, $v[\widehat{\theta }_{n}]\leq v[\tilde{%
\theta}_{n}]$ for all other unbiased estimators $\tilde{\theta}_{n}.$

In the case of unbiased estimators, $MSE\left( \hat{\theta}_{n},\theta
\right) =var\left( \hat{\theta}_{n}\right) ,$ and an efficient estimator
gives the smallest MSE\ compared to all other unbiased estimators.

Award: 5 marks for precise definition of efficiency, 5 marks for precise
relation to MSE, and 2.5 points for explanation, intuition.

(b) State the Weak Law of Large Numbers for a sequence of i.i.d. random
variables carefully stating any conditions required. Briefly explain the
theorem's importance for point estimation.

\bigskip

Solution:

\begin{theorem}[The Weak Law of Large Numbers]
If $X_{1},...,X_{n}$ are i.i.d. variables, and $\mathbb{E[}\left\vert
X_{1}\right\vert ]<\infty ,$ then $\frac{1}{n}\dsum\nolimits_{i=1}^{n}X_{i}%
\rightarrow _{p}\mu .$ Equivalently, $\left\vert \frac{1}{n}%
\dsum\nolimits_{i=1}^{n}X_{i}-\mathbb{E}[X_{1}]\right\vert \rightarrow
_{p}0. $
\end{theorem}

The WLLN tells us that the distribution of $\overline{X}_{n}$ piles up
around the true unknown $\mu .$

Award: 5 marks for statement + 4 marks for correct conditions; 3.5 marks for
explanation

(c) Define Type I and Type II errors in hypothesis testing and explain how
the probability of making such errors relates to the size and power of the
test.

\bigskip

Solution: A false rejection of the null hypothesis (rejecting $H_{0}$ when $%
H_{0}$ is true) is called Type I error. The probability of Type I error
(also referred to as size of the test) is given by $\mathbb{P}\left( \text{%
reject }H_{0}|H_{0}\text{ is true}\right) .$

A false acceptance of the null (accepting $H_{0}$ when $H_{0}$ is not true)
is called Type II error. The acceptance probability of the alternative is
called power of the test. The power is one minus probability of Type II
error: $\pi \left( \beta \right) =\mathbb{P}\left( \text{reject }H_{0}|H_{1}%
\text{ is true}\right) =\mathbb{P}\left( \left\vert z\right\vert
>z^{crit}|H_{1}\text{ is true}\right) $, which generally depends on the true
value of the parameter $\beta $, and for a well behaved test the power is
increasing both as $\beta $ moves away from the hypothesised value and when
the sample size increases.

\bigskip

Award: 2.5 marks for each precise definition of types of errors, and 2.5 for
each precise definition of size and power. Award the remaining 2.5 for
explanation.

(d) Suppose $\sqrt{n}\left( \hat{\theta}_{n}-\theta \right) \rightarrow _{d}%
\mathcal{N}\left( 0,1\right) $ as $n\rightarrow \infty .$ You want to test a
null hypothesis $\theta =0$ vesus a two-sided alternative $H_{1}:\theta \neq
0.$ Explain how to conduct a $z$ test at significance level $99\%$, and how
the $z$ statistic behaves under the null and alternative hypothesis.

\bigskip

Solution: The test statistic is $z=\frac{\widehat{\theta }_{j}}{1/\sqrt{n}}$

\QTR{frametitle}{U}nder the null ($\theta =0),$ $z\sim \mathcal{N}(0,1).$
Under the alternative $(\theta \neq 0),$ we have $z=\frac{\widehat{\theta }}{%
1/\sqrt{n}}=\frac{\left( \widehat{\theta }-\theta +\theta \right) }{1/\sqrt{n%
}}=\mathcal{N}(0,1)+\theta \sqrt{n}$ which diverges to $\pm $infinity as $n$
increases. We reject the null whenever $\left\vert z\right\vert >2.58$ at $%
99\%$ significance level.

Award: 5 point for correct expression of the statistic, 5 points for
chracterising behaviour under null and alternative (2.5 each); and 2.5 for
correct explanation of rejection decision. \bigskip

\textbf{Question 2}

Consider the linear regression model $y=X\beta +\varepsilon ,$ where $y$ is
an $n\times 1$ vector of dependent variables, $X$ is an $n\times k$ matrix
of regressors, $\beta $ is a $k\times 1$ vector of unknown parameters, and $%
\varepsilon $ is an $n\times 1$ vector of unobserved random variables. The
standard assumptions provided in the lecture are assumed to hold.

\begin{enumerate}
\item The linear regression model in is correctly specified and linear in $%
\beta $ and $\varepsilon .$

\item The columns of $X$ are linearly independent, so that $rank(X^{\prime
}X)=rank(X)=k<n.$

\item $\varepsilon _{i}$ are i.i.d. variables and $X$ are exogenous
identically distributed regressors, such that $\mathbb{E}[\varepsilon
_{i}|X]=0$ and $\mathbb{E}[\varepsilon _{i}^{2}|X]=\sigma ^{2}.$

\item $\limfunc{plim}_{n\rightarrow \infty }(\frac{1}{n}X^{\prime }X)=Q,$
where $\left\Vert Q\right\Vert <\infty $ and $Q$ is positive definite matrix.
\end{enumerate}

(a) Derive an expression for the variance of $\hat{\beta}^{OLS}$ under the
assumptions above$.$

\begin{gather*}
var[\widehat{\beta }_{n}]=\mathbb{E[(}\widehat{\beta }_{n}-\beta )(\widehat{%
\beta }_{n}-\beta )^{\prime }\mathbb{]} \\
=\mathbb{E}[\left( (X^{\prime }X)^{-1}X^{\prime }\varepsilon \right) \left(
(X^{\prime }X)^{-1}X^{\prime }\varepsilon \right) ^{\prime }] \\
=\mathbb{E}[(X^{\prime }X)^{-1}X^{\prime }\mathbb{E[}\varepsilon \varepsilon
^{\prime }|X]X(X^{\prime }X)^{-1}] \\
=\mathbb{E}[(X^{\prime }X)^{-1}X^{\prime }\sigma ^{2}I_{n}X(X^{\prime
}X)^{-1}]=\sigma ^{2}\mathbb{E}[(X^{\prime }X)^{-1}].
\end{gather*}

Award: 8.5 points for precise derivations and 4 for explaning steps.

(b) Show that $\hat{\beta}^{OLS}$ is consistent carefully explaining each
step of your derivations.

\bigskip

Solution:

$p\lim_{n\rightarrow \infty }\widehat{\beta }_{n}=\beta +p\lim_{n\rightarrow
\infty }\left[ \left( \frac{1}{n}\tsum\nolimits_{i=1}^{n}\mathbf{x}_{i}%
\mathbf{x}_{i}^{\prime }\right) ^{-1}\frac{1}{n}\tsum\nolimits_{i=1}^{n}%
\mathbf{x}_{i}\varepsilon _{i}\right] $%
\begin{eqnarray*}
&=&\beta +p\lim_{n\rightarrow \infty }\left( \frac{1}{n}\tsum%
\nolimits_{i=1}^{n}\mathbf{x}_{i}\mathbf{x}_{i}^{\prime }\right)
^{-1}p\lim_{n\rightarrow \infty }\left( \frac{1}{n}\tsum\nolimits_{i=1}^{n}%
\mathbf{x}_{i}\varepsilon _{i}\right) \text{ {\tiny by CMT}} \\
&=&\beta +\left( p\lim_{n\rightarrow \infty }\frac{1}{n}\tsum%
\nolimits_{i=1}^{n}\mathbf{x}_{i}\mathbf{x}_{i}^{\prime }\right)
^{-1}p\lim_{n\rightarrow \infty }\left( \frac{1}{n}\tsum\nolimits_{i=1}^{n}%
\mathbf{x}_{i}\varepsilon _{i}\right) \text{ {\tiny by CMT and Ass. 4}} \\
&=&\beta +Q^{-1}\mathbb{E[}\mathbf{x}_{i}\varepsilon _{i}]\text{ {\tiny by
CMT, WLLN and Ass. 4}} \\
&=&\beta \text{ \ }
\end{eqnarray*}%
since $\mathbb{E[}\mathbf{x}_{i}\varepsilon _{i}]=\mathbb{E[}\mathbf{x}_{i}%
\mathbb{E}\left( \varepsilon _{i}|x_{i}\right) ]=0$ by Assumption 3, note
that because of the i.i.d. assumptions, $\mathbb{E[}\mathbf{x}%
_{i}\varepsilon _{i}]=\mathbb{E[}\mathbf{x}_{1}\varepsilon _{1}].$ So, $%
\widehat{\beta }_{n}$ is a consistent estimator of $\beta .$

Award: 8.5 points for precise derivations and 4 for explaning steps.

(c) Suppose you know that $\beta =0,$ so that the model reduces to $%
y=\varepsilon .$ Is $\hat{\sigma}_{n}^{2}=\frac{1}{n}\varepsilon ^{\prime
}\varepsilon $ unbiased? Derive an expression for the variance of $\hat{%
\sigma}_{n}^{2}.$ 

Solution: $\hat{\sigma}^{2}=\frac{1}{n}\varepsilon ^{\prime }\varepsilon $
as the $\varepsilon $ is observed and $\mathbb{E}\left[ \hat{\sigma}^{2}%
\right] =\mathbb{E}\left[ \frac{1}{n}\tsum\nolimits_{i=1}^{n}\varepsilon
_{i}^{2}\right] =\sigma ^{2}.$ Alternatively, since $\varepsilon $ is
observed%
\begin{gather*}
\mathbb{E}[\widehat{\sigma }_{n}^{2}]=\mathbb{E}[\frac{1}{n}\widehat{%
\varepsilon }^{\prime }\widehat{\varepsilon }]=\mathbb{E}[\frac{1}{n}%
\varepsilon ^{\prime }\varepsilon ]=\frac{1}{n}\mathbb{E}[tr(\varepsilon
\varepsilon ^{\prime })] \\
=\frac{1}{n}tr(\mathbb{E}[\varepsilon \varepsilon ^{\prime }])=\frac{\sigma
^{2}}{n}\mathbb{E[}tr(I_{n})]=\sigma ^{2}\text{.}
\end{gather*}%
Let $z_{i}=\varepsilon _{i}^{2}-\sigma ^{2}$ which has expectation 0, 
\begin{gather*}
V\left( \widehat{\sigma }_{n}^{2}\right) =\mathbb{E}\left( \frac{1}{n}%
\tsum\nolimits_{i=1}^{n}\varepsilon _{i}^{2}-\sigma ^{2}\right) ^{2}= \\
=\mathbb{E}\left( \frac{1}{n}\tsum\nolimits_{i=1}^{n}\left( \varepsilon
_{i}^{2}-\sigma ^{2}\right) \right) ^{2} \\
=\frac{1}{n^{2}}\mathbb{E}\left( \tsum\nolimits_{i=1}^{n}z_{i}\right) ^{2}%
\overset{\text{0 mean, ind}}{=}\frac{1}{n^{2}}\tsum\nolimits_{i=1}^{n}%
\mathbb{E}z_{i}^{2} \\
\overset{\text{i.d.}}{=}\frac{1}{n}\left( \mathbb{E}z_{1}^{2}\right) =\frac{1%
}{n}\mathbb{E}\left( \varepsilon _{1}^{2}-\sigma ^{2}\right) ^{2}=\frac{1}{n}%
\left( \mathbb{E}\varepsilon _{1}^{4}-\sigma ^{4}\right) 
\end{gather*}
So we additionally require $\mathbb{E}\left( \varepsilon _{1}^{4}\right)
<\infty .$ $\hat{\sigma}_{n}^{2}$ is unbiased, this is because no additional
parameters need to be estimated. 

Award: 5 point for unbiasedness proof with precision, 5 points for the
variance derivation; 2.5 for explaining all steps, intuition, requiring 4th
moment.

\textbf{Question 3 }

(a) Consider the linear regression model from Question 2. Recall that the
ridge estimator is defined as $\hat{\beta}^{Ridge}=\left( X^{\prime
}X+\lambda _{n}I_{k}\right) ^{-1}X^{\prime }y.$ Suppose that $\lambda
_{n}=cn^{\delta /2}$ for some $\delta \geq 0.$ For which range of values for 
$c$ and $\delta $ is $\hat{\beta}^{Ridge}$ a consistent estimator for $\beta
?$

Solution: 
\begin{eqnarray*}
p\lim_{n\rightarrow \infty }\hat{\beta}^{Ridge} &=&\left(
p\lim_{n\rightarrow \infty }\frac{1}{n}\tsum\nolimits_{i=1}^{n}\mathbf{x}_{i}%
\mathbf{x}_{i}^{\prime }+\lim_{n\rightarrow \infty }cn^{\delta
/2-1}I_{k}\right) ^{-1}p\lim_{n\rightarrow \infty }\frac{1}{n}%
\tsum\nolimits_{i=1}^{n}\mathbf{x}_{i}\mathbf{x}_{i}^{\prime }\beta \\
&&+\left( p\lim_{n\rightarrow \infty }\frac{1}{n}\tsum\nolimits_{i=1}^{n}%
\mathbf{x}_{i}\mathbf{x}_{i}^{\prime }+\lim_{n\rightarrow \infty }cn^{\delta
/2-1}I_{k}\right) ^{-1}p\lim_{n\rightarrow \infty }\frac{1}{n}%
\tsum\nolimits_{i=1}^{n}\mathbf{x}_{i}\mathbf{\varepsilon }_{i}
\end{eqnarray*}

so for any values of $c\in \left( -\infty ,\infty \right) $ and any $\delta
\in \lbrack 0,2),$ $\lim_{n\rightarrow \infty }cn^{\delta /2-1}I_{k}=0$ so $%
p\lim_{n\rightarrow \infty }\hat{\beta}^{Ridge}=\beta .$

Award: 7.5 point for correct derivation of plim, 5 points for correct values
of $c$ and $\delta .$

\end{document}
